\subsection{An exponent bound when the weight is divisible by the residue prime}

\begin{corollary}[KT-ORDER-FIVE-PDIV-WEIGHT]
Let $p>3$ be prime and let $e\ge2$ be even, with
$\operatorname{ord}(\omega^e)=5$ and $p\mid e$.
For the strict group $\mathcal S_p(e)$ of (O5S.1), one has
\[
 \mathcal S_p(e)\simeq K_{2e}(\mathbf Z)\{p\}
       \simeq0\quad\hbox{or}\quad\mathbf Z/p\mathbf Z.
                                                        \tag{O5P.1}
\]
The displayed isomorphisms are isomorphisms of abstract finite groups.
This is an exponent bound in the stated weights; it does not prove
vanishing of the remaining order-five class component.
\end{corollary}

\begin{proof}
Put $L=\mathbf Q(\mu_p)$, $\Delta=\operatorname{Gal}(L/\mathbf Q)$,
$D=(\mathbf Q_p/\mathbf Z_p)(\chi_p^{-e})$, and $M_2=D[p^2]$.
The cyclotomic image of $G_L$ is $1+p\mathbf Z_p$.
Since $p\mid e$, the binomial congruence
$(1+pa)^e\equiv1\pmod {p^2}$ shows that $M_2$ is trivial on
$G_L$.  Its action thus factors through $\Delta$ and is
\[
                M_2\simeq(\mathbf Z/p^2\mathbf Z)(\widetilde\alpha),
                \qquad\alpha=\omega^{-e},                \tag{O5P.2}
\]
where $\widetilde\alpha$ is the Teichm\"uller lift of $\alpha$.
This lift is unique because $|\Delta|=p-1$ is prime to $p$.

The residual inertia character has order five.  An inertia element
on which $\alpha-1$ is nonzero therefore acts on $D$ with
$\chi_p^{-e}-1$ a unit.  Thus the global and $p$-local invariant
groups of $D$ vanish.  Applying cohomology to
\[
 0\longrightarrow M_2\longrightarrow D
       \xrightarrow{\ p^2\ }D\longrightarrow0             \tag{O5P.3}
\]
identifies $p^2$-torsion global classes with finite-coefficient
classes and preserves their strict local vanishing at $p$.
For $q\ne p$, inertia acts trivially on $D$, and its invariant
group $D$ is $p^2$-divisible.  The same exact sequence makes
$H^1(I_q,M_2)\longrightarrow H^1(I_q,D)$ injective, so the
finite-coefficient class is unramified there as well.

Restriction to $G_L$ is injective by inflation--restriction,
since $|\Delta|$ is prime to $p$.
It now gives an unramified homomorphism to $M_2$, and global
class field theory yields an injection
\[
 \mathcal S_p(e)[p^2]\lhook\joinrel\longrightarrow
 \operatorname{Hom}_{\mathbf Z_p[\Delta]}
   \bigl(A_p, (\mathbf Z/p^2\mathbf Z)(\widetilde\alpha)\bigr),
 \qquad A_p=\operatorname{Cl}(L)\{p\}.                   \tag{O5P.4}
\]
The equivariance has the same convention as (O5S.9).
Integral Teichm\"uller idempotents show that the right side
only uses $A_p^{[-e]}$.  KT-ORDER-FIVE-LENGTH-ONE gives this
component length at most one, so the right side has order
at most $p$.

By KT-ORDER-FIVE-STRICT-CYCLIC and the carried finiteness theorem,
$\mathcal S_p(e)\simeq\mathbf Z/p^a\mathbf Z$ for some $a\ge0$.
Its $p^2$-torsion has order $p^{\min(a,2)}$.
The bound from (O5P.4) therefore gives $a\le1$.
The full-group VK-PT comparison identifies the Pontryagin dual
of $K_{2e}(\mathbf Z)\{p\}$ with $\mathcal S_p(e)$,
which proves (O5P.1).
\end{proof}

\paragraph{Scope.}
The exponent conclusion uses the finite-level class-field injection
(O5P.4) and the full-group VK-PT comparison.
For $p\nmid e$, the coefficient $D[p^2]$ already has nontrivial
$G_L$-action, so this argument does not extend the exponent bound
to those weights.  The corollary leaves the same possible nonzero
class component; it removes no additional character or prime.
All class-number and foundational comparison inputs are the
separately checked theorems cited above, with their existing
source and review records.
